\section{Introduction}
Advanced Persistent Threats (APTs) are stealthy, multi-stage cyber campaigns that compromise enterprise systems over extended periods. Because APTs execute benign-looking commands and proceed through multi-hop lateral movements, traditional signature-based security systems are ineffective. Consequently, host audit logging and provenance graph analysis have emerged as powerful paradigms for APT detection. A provenance graph models system events (e.g., file reads, process spawns, network flows) as directed edges, preserving the causal history of system execution.

Despite their potential, existing state-of-the-art provenance-based IDS---such as OCR-APT \cite{ocrapt2025}, TraceCluster \cite{tracecluster2026}, and StageFinder \cite{stagefinder2026}---share a fatal limitation: \textbf{the clean-data assumption}. They require a pristine ``benign'' period or fully labeled training dataset to learn a statistical baseline of normal behavior. This assumption is unrealistic in practice for four reasons:
\begin{enumerate}
\item \textbf{Stealthy Contamination}: Attackers may already be active inside the enterprise environment when auditing begins, corrupting the ``normal'' training data.
\item \textbf{Concept Drift}: Enterprise systems naturally evolve (e.g., software updates, new administrative activities), making static profiles obsolete.
\item \textbf{Explosive Log Volumes}: Modern host environments generate millions of events daily, rendering manual labeling of audit logs intractable.
\item \textbf{False Positive Fatigue}: Traditional anomaly detection algorithms lack statistical controls on alarm rates, overwhelming security analysts.
\end{enumerate}

To address these challenges, we propose \textbf{ZeroCausal}, a framework that detects APT attacks using \textbf{uncurated streaming data with no manual labels, and no offline retraining}.

Our core insight is that \textbf{causal relationships are more stable than statistical correlations}. While statistical distributions of events vary over time due to concept drift, the underlying causal equations governing system execution remain invariant. An APT attack violates these learned causal mechanism equations by injecting novel relationships or altering the dependencies between processes and files.

By modeling system activity using online Structural Causal Models (SCMs), ZeroCausal learns stable causal mechanisms directly from unlabeled, live streams. Anomalies are quantified using a \textbf{Hybrid Causal Anomaly Score (H-CAS)} that fuses SCM residual errors and causal p-values into a single scalar. Crucially, to prevent false alarm fatigue, ZeroCausal incorporates a dynamic \textbf{conformal prediction feedback loop} that adjusts the detection threshold online, providing a user-defined False Positive Rate (FPR) budget.

\textbf{Relationship to Causal-IDS.}~The closest prior work is \textbf{Causal-IDS}~\cite{causalids2026}, which establishes the principle of detecting network intrusions as causal mechanism violations. ZeroCausal extends this insight to the provenance-graph domain and removes three deployment blockers that Causal-IDS cannot address: (1)~\emph{clean-data dependence}---Causal-IDS requires a pristine benign training corpus; ZeroCausal operates on raw, potentially contaminated streams; (2)~\emph{static offline operation}---Causal-IDS uses a fixed SCM with no drift recovery; ZeroCausal refits the SCM online upon detected concept drift; and (3)~\emph{uncontrolled false alarms}---Causal-IDS applies empirical fixed thresholds; ZeroCausal enforces a user-defined FPR budget via conformal prediction. H-CAS is also architecturally distinct from Causal-IDS's pure causal p-value score: the empirical residual CDF component in H-CAS provides a complementary signal that remains informative when individual causal p-values are imprecisely estimated under sparse or contaminated training data.

\subsection{Summary of Contributions}
\begin{itemize}
\item \textbf{Zero-Label Causal Graph IDS}: We design and implement ZeroCausal, one of the first APT detection systems in provenance graphs that combines online causal discovery and conformal false-alarm control to learn causal mechanisms online from raw, unlabeled, and potentially contaminated logs.
\item \textbf{Online Concept Drift Handling}: We integrate an \texttt{AdaptiveWindowDetector} that tracks multivariate streaming properties and refits SCM regression models online upon detecting structural changes.
\item \textbf{Conformal False Positive Control}: We utilize online conformal prediction with an adaptive feedback loop to automatically adjust the alert threshold to empirically control the user-defined target FPR budget.
\item \textbf{Significant Performance Optimizations}: By converting Pandas representations into raw 2D NumPy matrices, implementing a fast $O(\log N)$ binary search conformal check, and vectorizing empirical CDF calculations, we achieve a \textbf{13$\times$ streaming evaluation speedup} (reducing OpTC evaluation from 240 seconds to 18.88 seconds).
\item \textbf{Extensive Experimental Validation}: We benchmark ZeroCausal on five diverse datasets---three core benchmarks consisting of real enterprise host logs (OpTC) and two high-fidelity simulated environments (DARPA TC3 and NODLINK) designed to model known APT scenarios, plus two additional datasets (BETH and StreamSpot)---against Isolation Forest, LOF, One-Class SVM, and an unsupervised autoencoder, \emph{and} against these same baselines under contaminated training and concept drift. ZeroCausal achieves AUC of \textbf{0.9656 on BETH} (cross-platform Linux telemetry) and an average of \textbf{0.727 $\pm$ 0.071 on OpTC} over 10 seeds, outperforming Isolation Forest ($0.675 \pm 0.024$) while holding the empirical alarm rate to $4.1\%$ against a 5\% budget. We demonstrate that while ZeroCausal's causal-structure discovery remains stable under training contamination, its empirical-CDF calibration step is sensitive to contaminated calibration windows, causing overall performance to degrade, thus localizing the hardening target for future work. We also report an honest failure case on StreamSpot (AUC 0.4991), diagnosing limitations of linear SCMs on heterogeneous graph-type provenance data.
\end{itemize}
